Dans une compétition mathématique certains participants sont des amis. L'amitié est toujours réciproque. Un
groupe de participants est appelé une \emph{clique} si toute paire d'entre eux est formée de deux amis. (En
particulier, chaque groupe d'au plus un participant constitue une clique.) Le nombre de participants dans une
clique est appelé sa \emph{taille}.

On suppose que, dans cette compétition, la plus grande taille des cliques est paire. Montrer que les
participants peuvent être répartis dans deux pièces de telle sorte que la plus grande taille des cliques
contenues dans une de ces pièces soit égale à la plus grande taille des cliques contenues dans l'autre.
